\chapter[绝不是乐观主义]{绝不是乐观主义}

\chaptersubtitle{《不再依靠任何人》}

\chapternote{%
    ``Hope is \textbf{definitely not the same thing as optimism}.
    It is not the conviction that something will turn out well, but the certainty that something makes sense, regardless of how it turns out.''
    \par
    {\raggedleft\noindent ---Václav Havel, \emph{Disturbing the Peace}\par}
    \par
    “希望\textbf{绝不是乐观主义}。
    它不是相信某件事情会变得更好，而是确信，无论最终结果如何，这件事都有意义。” % 在本章中同步引用了该句子

    \par
    {\raggedleft\noindent ——瓦茨拉夫·哈维尔《扰乱和平》\par}
}

让我们重新回到起点。

除其他诸多意义外，《魔圆》也是对萌文化的批判。
所谓萌文化，是动画爱好者群体中的一种审美偏好，推崇柔弱、通常偏少女化、能唤起保护欲的角色。
然而，正如沙耶香篇中所揭示的那样，这种保护欲既是徒劳的，也是危险的。
因为一旦某人被置于“需要被保护”的位置，他或她就会被物化。
这样，此人便被视为一个需要被照料、被拯救的对象，而不再被承认为拥有自由能动性的完整主体。

正如魔法少女就是魔女一样，身为守护者也就是身为破坏者。
这一点既适用于沙耶香，也适用于焰（其姓“晓美\footnote{%
    译注：``whose family name, Akemi, means `fire', ...''，作者笔误。
    实际与“火焰”相应的是名“ほむら”，而非姓“晓美”。
}”被释为“火”，暗示其与萌文化中的“燃烧”有关）。
尤其值得注意的是，焰许下愿望时所追求的，并不是单纯地拯救小圆，而是与小圆交换位置：
她希望自己成为保护小圆的人，而不再是被小圆保护的人。
焰一旦把自己设定为小圆的守护者，也就同时把小圆固定在了“被守护者”的位置上。
正如下一集将进一步揭示，这种被守护的位置，正是小圆承受痛苦的原因。
而这并不难理解：
过度保护往往会使被守护者逐渐失去对自身能力和价值的自信。

动画中的佛教思想在此处尤为关键，因为焰正是被自己的观念所束缚。
她没有询问小圆真正想要什么，而是将自己所认定的“小圆应当想要什么”强加给她，这种判断归根结底来自焰自身的欲望。
因此，正如佛陀所言，欲望最终将焰导向了痛苦\cite{smart1998buddhism}。
同时，焰也误以为自己加诸世界的概念本身就是真实的。
可是，与此相对，世界本身便是无常且不断变化，焰自身的愿望和行动所带来的后果也处于其中。
这个变动不居的世界不断否定她试图建立的界限，例如未来与过去、生者与死者，以及人类、魔法少女和魔女之间的区分。
这也正契合禅宗的思想：
语言与思维会遮蔽人对整体真实的把握。
宇宙并非由彼此独立的实体和状态构成，而是一个始终处于变动不居的“一即一切”的整体\cite{smart1998formation-zen}。

焰最初试图把小圆从死亡中救出，因为她认为死亡不同于生命；
随后，她试图阻止小圆变成魔女，因为她认为魔女不同于魔法少女；
最后，她又试图阻止小圆成为魔法少女，因为她认为魔法少女不同于普通人类。
然而，焰从根本上并不明白自己身处怎样的故事——
值得注意的是，她在转入见泷原中学之前曾就读于一所基督教学校——
因此，她也无法认识到，自己的每一次行动都只会使局面进一步恶化。

简言之，这条路行不通。

让我们重新回到起点。

上一章中，我讨论了魔法少女类型的起源，以及“好女孩”与“坏女孩”（或“魔法少女”与“魔女”，或“圣母”与“荡妇”）这一二元结构如何限制女性力量的表达：
女性力量要么被纳入不会威胁男性霸权的形式之中，要么被置于外部者的位置上。
无论是森林深处的魔女，还是只能与他人无法看见的怪物战斗的魔法少女，都是这种外部者形象的不同表现。

第十集通过多条时间线塑造焰这一人物，并呈现她从非威胁性形象向外部者形象的转变。
在第一条时间线中，焰是一位无力而纯真的少女，脆弱到会成为魔女的受害者。
不过，这里也已经预示了她后来冷峻而坚定的一面：
尽管身体虚弱，她却是全动画唯一一位没有真正屈从于魔女精神支配的受害者。
随着时间线不断推进，焰逐渐变得更加强大，也越来越具有越界性。
第二条时间线开始时，她显然是所有魔法少女中最弱、也最缺乏战斗技巧的一位；
她的时间停止能力只是辅助能力，并没有明显的攻击用途，而她直接伤害魔女的能力也很有限。
然而，到这条时间线结束时，她已经能够将时间停止与自制炸弹结合起来，用以杀死魔女，而小圆则转而承担辅助角色。
到了第三条时间线，焰的力量与越界性进一步增强：
她不再只是制作炸弹，而是开始从黑道手中夺取枪械，从而具备了与人鱼魔女抗衡的能力。

在第四条时间线中，焰终于变成了我们熟悉的模样：
冷静、疏离、不带情感。
她既舍弃了少女式的双马尾和眼镜——
这两者都是“无害”的能指——
也舍弃了过去的犹豫不决，转而从作为霸权男性气质终极象征的军队那里夺取武器，开始独自战斗。

然而，焰的这些形象并不是按照从过去到未来的顺序依次展开，而是分散在不同的时间线之中。
也就是说，在某种意义上，这些不同阶段的焰同时存在。
供外部者战斗的“外部”空间，本身仍然是一种不会真正威胁男性霸权的渠道。
这两种形象之间看似存在的二元对立，和魔法少女与魔女之间的对立一样，都只是幻象；
麻美与焰不过是同一结构的两面。

因此，焰并不是对魔法少女类型的否定，也并没有真正挑战这一类型所依托的安全而无害的框架。
她只是这一框架的另一种表现形式，是另一种能被观众接受的魔法少女形象。

简言之，这条路行不通。

让我们重新回到起点。

在第一章中，我曾提出《魔圆》包含三个篇章，每一篇章都围绕某个角色或一组角色展开，并各自突出不同主题。
第九集收束了沙耶香与杏子篇。
这一篇章主要用于探讨抑郁，同时也使作品对魔法少女类型的反思形成完整回环。
因此，第十集便成为最终篇章的开端，而核心人物就是焰。

由此，作品安排了一整集来呈现焰的视角，展现她如何一次次回到过去，试图改变小圆的命运。
在这一过程中，焰似乎基本没有被那种曾在第二篇中压倒沙耶香的绝望支配。
本集中，焰唯一一次接近魔女化，是在第三条时间线里：
她提议与小圆一同变成魔女，并抹去整个世界和所有悲伤，这既预示了小圆最终许下的愿望，也预示了《叛逆》中的许多剧情。

第三条时间线的特殊之处在哪里？
第一条时间线中，焰是在小圆被瓦尔普吉斯之夜杀死之后才成为魔法少女的。
第二条时间线中，焰并未受伤，这表明她可能很少参与、甚至根本没有参与和瓦尔普吉斯之夜的战斗。
第四条时间线中，小圆则是在焰被瓦尔普吉斯之夜击败之后才成为魔法少女的。

也就是说，焰最接近魔女化的那条时间线，恰恰是她真正实现自己愿望的时间线：
她终于在与瓦尔普吉斯之夜的战斗中保护了小圆。
这正是沙耶香所说的希望与绝望的平衡；
希望的高峰会直通绝望的深渊。

焰之所以能够抵御绝望，是因为她并不以通常意义上的希望支撑自己。
她依靠的不是对美好结果的期待，而是自己的决心以及对小圆的爱。
由于这两者均非天然正面的情感，它们也就不会带来与之相抵的绝望。
换言之，焰体现了捷克剧作家、哲学家、革命者和政治家瓦茨拉夫·哈维尔对希望的理解：
“希望绝不是乐观主义。
它不是相信某件事情会变得更好，而是确信，无论最终结果如何，这件事都有意义。”
\cite{havel1990politics-hope}

正如焰在本集结尾所说，只要是为了小圆，无论要花多久时间，她都可以坚持下去。
那么，她在第八集中的崩溃又该如何理解？
小圆似乎在每一条时间线中都决意牺牲自己，而这正是焰的计划所面临的最大威胁。
焰实际上处于一种功能性抑郁\footnote{%
    译注：第八章曾提到，“功能性抑郁”是指个体虽然处于抑郁状态，却仍能在表面上维持日常生活、工作或社会功能，因此其痛苦不一定容易被外界察觉。
}的状态：
这意味着她虽然仍能行动，却同样处于抑郁之中；
只要再受到一点向下的推力，她就会崩溃。

简言之，这条路行不通。

让我们重新回到起点。

第十集呈现了一系列与主线故事相邻的时间线。
每当某一条时间线走向瓦尔普吉斯之夜带来的末日灾难时，焰都会将自己重置到离开医院的那一天，然后重新尝试。
然而，她似乎并不只是回到过去，而是在不同时间线之间移动；
因为即使不受焰介入的影响，某些要素也会在各条时间线中发生变化。
例如，在第一条时间线里，小圆在焰出院之前就已经成为了魔法少女。

在这四条不同时间线中，主时间线里的事件不断以回响——
或者说预兆——
的形式出现，这取决于观众如何理解。
从第一条时间线里小圆在学校走廊转身面对焰时给到鞋子的特写镜头（它呼应了第一集中同一场景里焰转身时类似的特写镜头），到杏子拒绝承认人鱼魔女已经完全取代了她所认识的沙耶香，这些细节都与主时间线中的事件彼此呼应。% TODO figure

更有意思的是那些细微的呼应。
比如，焰一次次在医院中醒来，呼应了第三集中夏洛特的魔女之种在医院孵化的情节，也进一步强化了二者之间的平行关系：
她们都是“真实剧集”的先兆，也都是麻美的敌手。
小圆恳求焰阻止自己与丘比签订契约，反而使焰代替小圆犯下了这一错误，呼应了第六集中询子的建议。
就连焰逐渐失去纯真，并转变为更加黑暗、更加成熟、也更加越界的人物这一过程，也与整个沙耶香篇形成呼应。

这些呼应不仅说明了不同时间线之间的关系，也对这些被呼应的事件作了再语境化。
从焰的视角来看，它们并不是事后才形成的呼应，而是先行的预示。
事实上，对观众来说，其中一些同样具有预示作用，例如麻美率先攻击其他魔法少女，便预示了《叛逆》中魔法少女彼此对立的局面。
这一切曾经发生过，而且很可能还会再次发生。

有两集彻底改变了整部作品。
初次观看时，这两集尤其突出，因为它们标志着作品从看似讲述某类故事到突然转变为另一种形态。
它们像镜像一样，恰好位于动画两端相对的位置：
第三集与倒数第三集。
本集就是其中的倒数第三集。

即使仅看结构（而事实上其他层面也相当丰富），《魔法少女小圆》也是一部极为严密的作品。
它的每一个瞬间都经过精心安排，在短得惊人的篇幅内推进复杂的故事，并充分塑造人物。
一个例子便是贯穿全剧的对称性：
许多或大或小的情节都会在恰当的位置再次出现。

简言之，这条路行不通。
